\section{Mathematical model}
\label{sec:mathematical_model}
In this chapter we introduce the integer linear program used to solve the instances. 
The model is formulated as a two-stage stochastic mixed-integer linear program and
is solved through its deterministic equivalent formulation.

\subsection{Sets and Indices}

\begin{itemize}
  \item $p \in P$ -- set of patients (surgeries)
  \item $P_0 = P \cup \{0\}$ -- patients plus dummy depot node (session start)
  \item $h \in H$ -- set of available OR sessions
  \item $s \in S$ -- set of scenarios
  \item $q \in Q = \{\text{CHI},\text{ORT},\text{KNO}\}$ -- set of surgical
        specialties
\end{itemize}

\subsection{Parameters}

\begin{itemize}
  \item $d_{ps}$ -- surgery duration of patient $p$ in scenario $s$
  \item $\pi_s$ -- probability of scenario $s$,
        where $\sum_{s \in S} \pi_s = 1$
  \item $c$ -- changeover time
  \item $t^{\text{start}}$ -- common session start time (identical for all ORs)
  \item $t^{\text{close}}$ -- common session closing time (identical for all
        ORs)
  \item $q_{pq} \in \{0,1\}$ -- 1 if patient $p$ belongs to specialty $q$
  \item $M$ -- sufficiently large constant (big-$M$)
  \item $\beta_1=0.6, \beta_2=0.2, \beta_3=0.2, \beta_4=100$ -- weights for waiting time, idle
        time, overtime, and specialty mismatch, respectively
\end{itemize}

\subsection{First-Stage Decision Variables}

The following decisions are made before scenario realisation:

\begin{itemize}
  \item $A_p$ -- appointment time assigned to patient $p$
  \item $X_{ijh}$ -- 1 if patient $j$ is scheduled directly after
        patient $i$ in session $h$
  \item $Y_{ph}$ -- 1 if patient $p$ is assigned to session $h$
  \item $Z_h$ -- 1 if OR session $h$ is opened
  \item $V_{hq}$ -- 1 if session $h$ is assigned specialty $q$
  \item $D_{hq}$ -- number of patients in session $h$ with specialty $q$
        that are not placed in a session with their specialty
  \item $U_{ph}$ -- position of patient $p$ in session
        $h$ (MTZ auxiliary variable)
\end{itemize}

\subsection{Second-Stage Decision Variables}

The following variables are evaluated per scenario $s$:

\begin{itemize}
  \item $S_{ps}$ -- start time of patient $p$ under
        scenario $s$
  \item $W_{ps}$ -- waiting time of patient $p$ under scenario $s$
  \item $I_{pp's}$ -- idle time in between patients $p$ and $p'$ under scenario $s$
  \item $O_{hs}$ -- overtime of session $h$ under scenario $s$
\end{itemize}

\subsection{Two-Stage Stochastic Formulation}

The model has a two-stage stochastic structure. First-stage decisions are made
before surgery durations are known. These decisions include appointment times,
session opening decisions, patient-to-session assignments, sequencing decisions,
specialty assignments, violation-counting variables, and MTZ auxiliary variables.

Let

\[
x = (A_p, X_{ijh}, Y_{ph}, Z_h, V_{hq}, D_{hq}, U_{ph})
\]

denote the vector of first-stage variables.

After a scenario $s \in S$ is realised, the second-stage variables determine
the realised start times, waiting times, idle times, and overtime. Let

\[
y_s = (S_{ps}, W_{ps}, I_{pp's}, O_{hs})
\]

denote the vector of second-stage variables in scenario $s$.

For a fixed first-stage solution $x$ and realised scenario $s$, the second-stage
problem is formulated as the recourse function $Q(x,s)$, which gives the minimum
operational cost achievable under scenario $s$ with surgery durations $d_{ps}$.
The first-stage solution $x$ enters $Q(x,s)$ entirely as fixed parameters:
$A_p$ fixes each patient's appointment time, while $X_{pp'h}$ and $Y_{ph}$
fix the sequencing and session assignments respectively.

\begin{align}
  Q(x, s) = \min_{S,W,I,O} \quad
    & \beta_1 \sum_{p \in P} W_{ps}
    + \beta_2 \sum_{p \in P}\sum_{\substack{p' \in P \\ p' \neq p}} I_{pp's}
    + \beta_3 \sum_{h \in H} O_{hs}
  \label{eq:recourse_obj}
\end{align}

\noindent subject to the following scenario-dependent constraints.

\paragraph{Start time propagation.}
If patient $p'$ follows patient $p$ in session $h$, the start of $p'$ must
respect patient $p$'s realised duration and the changeover time $c$:

\begin{align}
  S_{p's} \geq\; & S_{ps} + d_{ps} + c - M\bigl(1 - \sum_{h \in H} X_{pp'h}\bigr)
  \nonumber \\
                 & \forall\; p, p' \in P,\; p \neq p',\; s \in S
  \label{eq:rc_seq}
\end{align}

\paragraph{Appointment lower bound.}
A patient cannot start surgery before their appointed time:

\begin{align}
  S_{ps} \geq\; A_p
  \qquad \forall\; p \in P,\; s \in S
  \label{eq:rc_appt}
\end{align}

\paragraph{Waiting time.}
Waiting time is the excess of actual start time over appointment time:

\begin{equation}
  W_{ps} \geq S_{ps} - A_p
  \qquad \forall\; p \in P,\; s \in S
  \label{eq:rc_wait}
\end{equation}

\paragraph{Overtime.}
Overtime is incurred when any patient's finish time exceeds the session closing
time. The big-$M$ term deactivates the constraint for patients not in session
$h$:

\begin{align}
  O_{hs} \geq\; S_{ps} + d_{ps} - t^{\text{close}} - M(1 - Y_{ph})
  \qquad \forall\; p \in P,\; h \in H,\; s \in S
  \label{eq:rc_ot}
\end{align}

\paragraph{Idle time.}
Idle time between two consecutively scheduled patients $p$ and $p'$ is the
gap between the finish time of $p$ and the start time of $p'$, minus the
mandatory changeover time. The big-$M$ term deactivates the constraint for
non-consecutive pairs:

\begin{align}
  I_{pp's} \geq\; & S_{p's} - S_{ps} - d_{ps} - c
                  - M\Bigl(1 - \sum_{h \in H} X_{pp'h}\Bigr)
  \nonumber \\
  & \forall\; p, p' \in P,\; p \neq p',\; s \in S
  \label{eq:rc_idle}
\end{align}

\paragraph{Non-negativity.}

\begin{equation}
  S_{ps},\; W_{ps},\; I_{pp's},\; O_{hs} \geq 0
  \qquad \forall\; p,p' \in P,\; p \neq p',\; h \in H,\; s \in S
  \label{eq:rc_nn}
\end{equation}

\noindent The first-stage variables $A_p$, $X_{ijh}$, $Y_{ph}$, $Z_h$,
$V_{hq}$, $D_{hq}$, and $U_{ph}$ enter $Q(x,s)$ as fixed parameters. Given
these, all constraints are linear in the second-stage variables, so $Q(x, s)$
is a linear program for each $(x, s)$.

The full two-stage stochastic problem can therefore be written as:

\begin{align}
  \min_x \quad
    & \beta_4 \sum_{h \in H} \sum_{q \in Q} D_{hq}
    + \sum_{s \in S} \pi_s Q(x, s)
  \label{eq:two_stage_obj}
\end{align}

where the first term represents the scenario-independent specialty mismatch
penalty, and the second term represents the expected second-stage operational
cost.

\subsection{Deterministic Equivalent MILP}

Rather than solving $Q(x,s)$ separately for each scenario, the two-stage
model is reformulated as a single deterministic equivalent by replicating
the second-stage variables and constraints across all scenarios $s \in S$.
The deterministic equivalent is a mixed-integer linear program (MILP) formulated in VRP-style.

\subsubsection{Expanded Objective Function}

The objective minimises the expected weighted sum of patient waiting time, OR
idle time, and OR overtime, together with a scenario-independent penalty for
specialty violations:

\begin{align}
  \min \quad
    & \sum_{s \in S} \pi_s \Bigl[
        \beta_1 \sum_{p \in P} W_{ps}
      + \beta_2 \sum_{p \in P}\sum_{\substack{p' \in P \\ p' \neq p}} I_{pp's}
      + \beta_3 \sum_{h \in H} O_{hs}
      \Bigr] \nonumber \\
    & + \; \beta_4 \sum_{h \in H} \sum_{q \in Q} D_{hq}
  \label{eq:obj}
\end{align}

The mismatch penalty is scenario-independent because specialty and assignment
decisions are first-stage.

\subsubsection{First-Stage Constraints}

\paragraph{Appointment time bounds.}

Each appointment must fall within the OR session window, and a patient scheduled
first within a session should be scheduled at the opening time of the session.

\begin{equation}
  t^{\text{start}} \leq A_p \leq t^{\text{close}}
  \qquad \forall\, p \in P
  \label{eq:appt_bounds}
\end{equation}

\begin{equation}
  A_p \leq t^{\text{start}}+M\left(1-\sum_{h \in H}X_{0ph}\right)
  \qquad \forall\, p \in P
  \label{eq:start_first}
\end{equation}

\paragraph{Assignment and session opening.}

Each patient must be assigned to exactly one session:

\begin{equation}
  \sum_{h \in H} Y_{ph} = 1
  \qquad \forall\, p \in P
  \label{eq:assign}
\end{equation}

A patient can only be assigned to an opened session:

\begin{equation}
  Y_{ph} \leq Z_h
  \qquad \forall\, p \in P,\; h \in H
  \label{eq:open}
\end{equation}

\paragraph{Specialty assignment.}

Each opened session is assigned exactly one specialty; closed sessions receive
none:

\begin{equation}
  \sum_{q \in Q} V_{hq} = Z_h
  \qquad \forall\, h \in H
  \label{eq:spec}
\end{equation}

\paragraph{Specialty violation counting.}

The integer variable $D_{hq}$ counts the number of patients in session $h$
whose specialty does not match $q$. For each session $h$ and specialty $q$,
the number of assigned patients belonging to specialty $q$ cannot exceed the
session capacity unless a violation is recorded:

\begin{equation}
  \sum_{p \in P} Y_{ph} \cdot q_{pq}
  \leq M \cdot V_{hq} + D_{hq}
  \qquad \forall\, h \in H,\; q \in Q
  \label{eq:viol}
\end{equation}

When session $h$ is assigned specialty $q$ (i.e.\ $V_{hq} = 1$), the big-$M$
term is active and $D_{hq} = 0$ is optimal under minimisation provided all
assigned patients match. When $V_{hq} = 0$, the big-$M$ term vanishes and
$D_{hq}$ must be at least as large as the number of patients of specialty $q$
placed in session $h$, recording each as a violation.
This is a \emph{soft} penalty rather than a hard constraint: specialty mixing is
discouraged through the weighted term $\beta_4 \sum_{h,q} D_{hq}$ in the
objective, which matches the data observation that KNO can share a session with
other specialties but that this should be avoided where possible. A large
$\beta_4 = 100$ makes such mixing strongly unattractive while keeping the model
feasible when an unavoidable mismatch occurs.

\paragraph{Routing --- flow conservation.}

The sequencing variables $X_{ijh}$ define a single route per open session,
mirroring a vehicle routing formulation:

\begin{align}
  \sum_{p \in P} X_{0ph} &= Z_h
  \qquad \forall\, h \in H
  \label{eq:depot} \\[2pt]
  \sum_{\substack{p' \in P_0 \\ p' \neq p}} X_{pp'h} &= Y_{ph}
  \qquad \forall\, p \in P,\; h \in H
  \label{eq:flow_out} \\[2pt]
  \sum_{\substack{p \in P_0 \\ p \neq p'}} X_{pp'h} &= Y_{p'h}
  \qquad \forall\, p' \in P,\; h \in H
  \label{eq:flow_in}
\end{align}

\paragraph{Sub-tour elimination.}

Miller--Tucker--Zemlin (MTZ) constraints prevent disconnected cycles within a
session:

\begin{align}
  U_{ph} - U_{p'h} + |P|\, X_{pp'h} \leq |P| - 1 \nonumber \\
  \forall\, p,p' \in P,\; p \neq p',\; h \in H
  \label{eq:mtz}
\end{align}

\paragraph{First-stage domains.}

\begin{align}
  & X_{ijh},\; Y_{ph},\; Z_h,\; V_{hq} \in \{0,1\}
  \label{eq:dom_bin} \\[2pt]
  & D_{hq},\; U_{ph} \in \mathbb{Z}_{\geq 0}
  \qquad
  \label{eq:dom_viol} \\[2pt]
  & A_p \geq 0
  \label{eq:dom_appt}
\end{align}

\subsubsection{Scenario-Dependent Second-Stage Constraints}

The following constraints are included for each scenario $s \in S$.

\paragraph{Timing --- start time propagation.}

If patient $p'$ immediately follows patient $p$ in session $h$, the start time
of $p'$ must respect $p$'s realised duration and the changeover time:

\begin{align}
  S_{p's} \geq\; & S_{ps} + d_{ps} + c - M\bigl(1 - \sum_{h \in H} X_{pp'h}\bigr)
  \nonumber \\
                 & \forall\; p, p' \in P,\; p \neq p', s \in S
  \label{eq:time_seq}
\end{align}

\paragraph{Session lower bounds.}

A patient's surgery cannot start before the session opens or before the
patient's appointment time:

\begin{align}
  S_{ps} &\geq A_p
  \qquad \forall\, p \in P,\; s \in S
  \label{eq:lb_appt}
\end{align}

\paragraph{Waiting time.}

Patient waiting time is the positive difference between the actual start time
and the planned start time, i.e., the appointment time:

\begin{equation}
  W_{ps} \geq S_{ps} - A_p
  \qquad \forall\, p \in P,\; s \in S
  \label{eq:wait}
\end{equation}

\paragraph{Overtime.}

Overtime in session $h$ under scenario $s$ is the amount of time by which the
last patient's finish time exceeds the closing time:

\begin{align}
  O_{hs} \geq\; &
                    S_{ps} + d_{ps}
                  - t^{\text{close}}-M(1- Y_{ph}) \nonumber \\
                & \forall\, p \in P,\;\forall\, h \in H,\; s \in S
  \label{eq:overtime}
\end{align}

\paragraph{Idle time.}

The idle time is defined as the time between the end time and starting time of
two consecutively scheduled surgeries minus the mandatory changeover time
between two subsequent surgeries:

\begin{align}
  I_{pp's} \geq\; & S_{p's}-S_{ps}-d_{ps}-
                  c-M\left(1-\sum_{h \in H}X_{pp'h}\right) \nonumber \\
                &  \forall\, p,p' \in P,\; p \neq p',\; s \in S
  \label{eq:idle}
\end{align}

\paragraph{Second-stage domains.}

\begin{equation}
  S_{ps},\; W_{ps},\; I_{pp's},\; O_{hs} \geq 0
  \qquad \forall\, p,p' \in P,\; p \neq p',\; h \in H,\; s \in S
  \label{eq:dom_second}
\end{equation}

The deterministic equivalent is obtained by combining the expanded objective
\eqref{eq:obj}, the first-stage constraints
\eqref{eq:appt_bounds}--\eqref{eq:dom_appt}, and the scenario-dependent
second-stage constraints \eqref{eq:time_seq}--\eqref{eq:dom_second} for all
scenarios $s \in S$.

The DE is a mixed-integer linear program (MILP): all integrality requirements
appear in the first stage, and the second-stage subproblem is a pure LP for
any fixed first-stage solution. The total number of second-stage variables
scales as $\mathcal{O}(|S| \cdot (|P| + |H|))$, so tractability depends on
keeping $|S|$ moderate, for instance through Sample Average Approximation
(SAA), in which $|S|$ scenarios are drawn from the duration distribution and
the DE is solved once.

\subsection{Experimental Model Variants}

\subsubsection{Step 1: Timing Only}

In the first step we only evaluate the timing of treatments already being
assigned to a session and a position within the session by constraining
$X_{ijh}$.

\subsubsection{Step 2: Sequencing + Timing}

In the second step we evaluate the timing and sequencing of treatments that are
already assigned to a session, by constraining $Y_{ph}$.

\subsubsection{Step 3: Full Model}

In the third step we evaluate the timing, sequencing, and assignment of
treatments by evaluating the full model.

\paragraph{Computational tractability.}
The full model is a large-scale MILP whose LP relaxation is weakened by the
Big-$M$ coefficients in equations~2, 5, 6, 11, 15, 23, 26 and 27.
To reduce the solution space, two families of symmetry-breaking constraints are
added when the session assignment is free (Step~3 only).
\textbf{SB1} enforces a canonical specialty ordering across sessions: sessions
of specialty~$q_1$ must occupy lower session indices than sessions of
specialty~$q_2$ whenever $q_1$ precedes $q_2$ in the ordering of $Q$.
\textbf{SB2} breaks within-specialty symmetry by requiring the first patient in
a lower-indexed session to have a strictly lower patient ID than the first
patient in any higher-indexed session of the same specialty.
Additionally, Step~3 is warm-started from the Step~2 solution via a
within-specialty session-permutation that satisfies SB1 and SB2, injecting a
high-quality incumbent before branching begins.
Despite these measures, a 600-second time limit and a $5\%$ MIP-gap tolerance
are applied in SAA replications to balance solution quality with run time.